% TODO make hw stylesheet
\documentclass[11pt]{scrartcl}
\usepackage[a4paper, total={6in, 8in}]{geometry}
\usepackage[usenames,dvipsnames,svgnames]{xcolor}
\usepackage[shortlabels]{enumitem}
\usepackage[framemethod=TikZ]{mdframed}
\usepackage{amsmath,amssymb,amsthm}
\usepackage{accents}
\usepackage[colorlinks]{hyperref}
\usepackage{multicol}
\usepackage{tabularx}
\usepackage{stmaryrd}

\hypersetup{
	colorlinks=true,
	linkcolor=blue,
	filecolor=magenta,
	urlcolor=magenta,
	pdftitle={MAT 528 Notes}
}

\usepackage{microtype}
\usepackage{mathtools}
\usepackage[headsepline]{scrlayer-scrpage}
\usepackage{thmtools}
\usepackage[textsize=footnotesize]{todonotes}

\mdfdefinestyle{mdbluebox}{roundcorner=10pt,innerbottommargin=9pt,
	linecolor=blue,backgroundcolor=TealBlue!5,}
\declaretheoremstyle[headfont=\sffamily\bfseries\color{MidnightBlue},
mdframed={style=mdbluebox},]{thmbluebox}

\mdfdefinestyle{mdbloobox}{frametitlefont=\bfseries,innerbottommargin=8pt,
	nobreak=true,backgroundcolor=TealBlue!5,linecolor=blue,}
\declaretheoremstyle[headfont=\bfseries\color{MidnightBlue},
mdframed={style=mdbloobox},headpunct={\\[3pt]},postheadspace=0pt,]{thmbloobox}

\mdfdefinestyle{mdredbox}{frametitlefont=\bfseries,innerbottommargin=8pt,
	nobreak=true,backgroundcolor=Salmon!5,linecolor=RawSienna,}
\declaretheoremstyle[headfont=\bfseries\color{RawSienna},
mdframed={style=mdredbox},headpunct={\\[3pt]},postheadspace=0pt,]{thmredbox}

\mdfdefinestyle{mdgreenbox}{linecolor=ForestGreen,backgroundcolor=ForestGreen!5,
	linewidth=2pt,rightline=false,leftline=true,topline=false,bottomline=false,}
\declaretheoremstyle[headfont=\bfseries\sffamily\color{ForestGreen!70!black},
mdframed={style=mdgreenbox},headpunct={ --- },]{thmgreenbox}

\mdfdefinestyle{mdorangebox}{linecolor=RedOrange,backgroundcolor=RedOrange!5,
	linewidth=2pt,rightline=false,leftline=true,topline=false,bottomline=false,}
\declaretheoremstyle[headfont=\bfseries\sffamily\color{RedOrange!70!black},
mdframed={style=mdorangebox},headpunct={ --- },]{thmorangebox}

\mdfdefinestyle{mdblackbox}{linecolor=black,backgroundcolor=RedViolet!5!gray!5,
	linewidth=3pt,rightline=false,leftline=true,topline=false,bottomline=false,}
\declaretheoremstyle[mdframed={style=mdblackbox}]{thmblackbox}

\declaretheoremstyle[spaceabove=6pt,spacebelow=6pt]{basehead}

\declaretheorem[style=thmbluebox,name=Theorem,numberwithin=section]{theorem}
\declaretheorem[style=thmredbox,name=Example,sibling=theorem]{example}
\declaretheorem[style=thmbluebox,name=Corollary,sibling=theorem]{corollary}
\declaretheorem[style=thmbluebox,name=Proposition,sibling=theorem]{prop}
\declaretheorem[style=thmbluebox,name=Lemma,sibling=theorem]{lemma}
\declaretheorem[style=thmbluebox,name=Lemma,numberwithin=theorem]{thmlemma}
\declaretheorem[style=thmbluebox,name=Theorem,numbered=no]{theorem*}
\declaretheorem[style=thmbluebox,name=Lemma,numbered=no]{lemma*}
\declaretheorem[style=thmgreenbox,name=Claim,numberwithin=theorem]{claim}
\declaretheorem[style=thmgreenbox,name=Claim,numbered=no]{claim*}
\declaretheorem[style=thmblackbox,name=Remark,numberwithin=theorem]{remark}
\declaretheorem[style=thmblackbox,name=Remark,numbered=no]{remark*}
\declaretheorem[style=thmgreenbox,name=Definition,sibling=theorem]{definition}
\declaretheorem[style=thmgreenbox,name=Definition,numbered=no]{definition*}
\declaretheorem[style=thmorangebox,name=Warning,sibling=theorem]{warning}
\declaretheorem[style=thmorangebox,name=Warning,numbered=no]{warning*}
\declaretheorem[style=thmblackbox,name=Exercise,sibling=theorem]{exercise}
\declaretheorem[style=thmblackbox,name=Exercise,numbered=no]{exercise*}

\addtolength{\textheight}{3.14cm}
\providecommand{\ol}{\overline}
\providecommand{\ceil}[1]{\left\lceil #1 \right\rceil}
\providecommand{\floor}[1]{\left\lfloor #1 \right\rfloor}
\newcommand{\dd}{\mathrm{d}}
\providecommand{\interior}[1]{\accentset{\circ}{#1}}
\providecommand{\ul}{\underline}
\providecommand{\eps}{\varepsilon}
\providecommand{\half}{\frac{1}{2}}
\providecommand{\dang}{\measuredangle} %% Directed angle
\providecommand{\CC}{\mathbb C}
\providecommand{\FF}{\mathbb F}
\providecommand{\NN}{\mathbb N}
\providecommand{\QQ}{\mathbb Q}
\providecommand{\RR}{\mathbb R}
\providecommand{\ZZ}{\mathbb Z}
\providecommand{\dg}{^\circ}
\providecommand{\ii}{\item}
\providecommand{\setdiff}{\smallsetminus}
\providecommand{\alert}[1]{{\sffamily\textbf{\textcolor{blue}{#1}}}}
\providecommand{\inv}{^{-1}}
\providecommand{\mb}[1]{\mathbf{#1}}

\renewcommand{\tilde}{\widetilde}

\reversemarginpar

\newenvironment{soln}{\begin{proof}[Solution]}{\end{proof}}
\newenvironment{walkthrough}{\noindent \textbf{Walkthrough.}}{\medskip}
\newlist{walk}{enumerate}{3}
\setlist[walk]{label=\bfseries (\alph*)}

\providecommand{\pow}{\operatorname{Pow}}
\providecommand{\vocab}[1]{{\textbf{\textcolor{ForestGreen}{#1}}}}
\providecommand{\lcm}{\operatorname{lcm}}
\providecommand{\abs}[1]{\left|#1\right|}

\usepackage{asymptote}
\begin{asydef}
	defaultpen(fontsize(10pt));
	unitsize(1cm);
	size(8cm); // set a reasonable default
	usepackage("amsmath");
	usepackage("amssymb");
	settings.tex="pdflatex";
	settings.outformat="pdf";
	// Replacement for olympiad+cse5 which is not standard
	import geometry;
	// recalibrate fill and filldraw for conics
	void filldraw(picture pic = currentpicture, conic g, pen fillpen=defaultpen, pen drawpen=defaultpen)
	{ filldraw(pic, (path) g, fillpen, drawpen); }
	void fill(picture pic = currentpicture, conic g, pen p=defaultpen)
	{ filldraw(pic, (path) g, p); }
	// some geometry
	pair foot(pair P, pair A, pair B) { return foot(triangle(A,B,P).VC); }
	pair orthocenter(pair A, pair B, pair C) { return orthocentercenter(A,B,C); }
	pair centroid(pair A, pair B, pair C) { return (A+B+C)/3; }
	// cse5 abbreviations
	path CP(pair P, pair A) { return circle(P, abs(A-P)); }
	path CR(pair P, real r) { return circle(P, r); }
	pair IP(path p, path q) { return intersectionpoints(p,q)[0]; }
	pair OP(path p, path q) { return intersectionpoints(p,q)[1]; }
	path Line(pair A, pair B, real a=0.6, real b=a) { return (a*(A-B)+A)--(b*(B-A)+B); }
	// cse5 more useful functions
	picture CC() {
		picture p=rotate(0)*currentpicture;
		currentpicture.erase();
		return p;
	}
	pair MP(Label s, pair A, pair B = plain.S, pen p = defaultpen) {
		Label L = s;
		L.s = "$"+s.s+"$";
		label(L, A, B, p);
		return A;
	}
	pair Drawing(Label s = "", pair A, pair B = plain.S, pen p = defaultpen) {
		dot(MP(s, A, B, p), p);
		return A;
	}
	path Drawing(path g, pen p = defaultpen, arrowbar ar = None) {
		draw(g, p, ar);
		return g;
	}
\end{asydef}

\usepackage{mathrsfs}

\ihead{\footnotesize MAT 528 Notes}
\ohead{\footnotesize \today}

\title{\vspace{-2em}MAT 528: Foundations of Topology}
\author{\large Aditya Pahuja}
\date{}
\begin{document}
\maketitle
\tableofcontents
\pagebreak

\section{August 25}\label{sec:aug25}
\begin{definition}[Topological space]
    \label{def:topspace}
    Given a set $X$, a topology $\mathcal T$ on $X$ is
    a collection of subsets of $X$, called \vocab{open sets}, such that
    \begin{enumerate}[(1)]
	\ii $\varnothing$ and $X$ are in $\mathcal T$.
	\ii If $\{U_\alpha\}_{\alpha\in I}$ is
	an arbitrary subcollection of $\mathcal T$, then
	$\bigcup_{\alpha\in I}U_\alpha$ is also in $\mathcal T$.
	\ii If $U_1,\dots,U_n$ are finitely many sets in $\mathcal T$,
	then $U_1\cap\cdots\cap U_n$ is in $\mathcal T$ as well.
    \end{enumerate}
    In English, (2) says that arbitrary unions of open sets are also open sets
    and (3) says that finite intersections of open sets are also open sets.
    The pair $(X,\mathcal T)$ is a \vocab{topological space},
    which, if the underlying topology is clear from context,
    will just be referred to as $X$.
\end{definition}

\begin{example}
    Suppose $X = \{a,b,c\}$ is a $3$-element set.
    Here are some topologies on $X$:
    \begin{enumerate}[(a)]
	\ii $\mathcal T = \{\varnothing, X\}$;
	in general, this is called the \vocab{trivial topology}.
	\ii $\mathcal T = \mathcal P(X)$, the power set of $X$;
	in general, this is called the \vocab{discrete topology}.
	\ii $\mathcal T = \{\{a\}, \varnothing, X\}$.
	\ii $\mathcal T = \{\{a\}, \{a,b\}, \varnothing, X\}$.
    \end{enumerate}
    Some non-examples: $\{\{a\}, \{b\}, \varnothing, X\}$
    and $\{\{a,b\}, \{a,c\}, \varnothing, X\}$ aren't topologies on $X$.
\end{example}

A very standard example that lends itself well to visual intuition is the following:
\begin{example}
    [Euclidean topology on $\RR$]
    \label{ex:euclideantopology-R}
    Give $\RR$ a topology by defining $U\subseteq\RR$ to be open
    if for all $x\in U$, there is $\eps > 0$ such that $(x-\eps,x+\eps)\subseteq\RR$.
    Checking that this is a topology:
    \begin{enumerate}[(1)]
        \ii $\varnothing$ and $X$ are trivially open by this definition (check it!).
	\ii If $\{U_\alpha\}_{\alpha\in I}$ is an arbitrary collection of
	open subsets of $\RR$, then any $x\in\bigcup_{\alpha\in I}U_\alpha$
	is contained in one of the $U_\alpha$s, say $U_\beta$;
	then, there exists an $\eps>0$ such that
	$(x-\eps,x+\eps)\subseteq U_\beta\subseteq\bigcup_{\alpha\in I} U_\alpha$
	so the union of the $U_\alpha$ is open.
	\ii If $U_1$, \dots, $U_n$ are open and $x$ lies in their intersection,
	then it of course lies in each $U_i$. Thus for each $i=1,\dots,n$,
	there is $\eps_i>0$ such that $(x-\eps_i,x+\eps_i)\subseteq U_i$.
	Therefore, letting $\eps$ be the positive number $\min_{i=1,\dots,n}\eps_i$,
	we see that $U_1\cap\cdots\cap U_n$ is open because
	$(x-\eps,x+\eps)\subseteq (x-\eps_i,x+\eps_i)\subseteq U_i$ for each $i$.
    \end{enumerate}
\end{example}

\begin{exercise}
    \label{exer:extrapolate-euclideantop}
    Extrapolate the construction of Example \ref{ex:euclideantopology-R}
    to get a topology on $\RR^n$ (there are a few natural ways to do this;
    try to show that they all result in the same topology).
\end{exercise}

\begin{definition}[Basis]
    \label{def:basis}
    Let $\mathcal T$ be a topology of $X$.
    Then, $\mathcal B\subseteq\mathcal T$ is
    a \vocab{basis} for $\mathcal T$ if any $U\subseteq\mathcal T$
    can be written as the union of some elements of $\mathcal B$.
\end{definition}

\begin{example}
    [Basis of the Euclidean topology]
    \label{ex:euclideantopology-basis}
    Consider $\RR$ with the Euclidean topology.
    Then, $\mathcal B = \{(a,b) : a,b\in\RR,\; a<b\}$
    is a basis because if $U$ is an open subset of $\RR$,
    then for all $x\in U$, there is an $\eps_x>0$ such that
    $(x-\eps_x,x+\eps_x)\subseteq U$, which implies that
    \begin{equation*}
	\bigcup_{x\in U} (x-\eps_x,x+\eps_x) = U.
    \end{equation*}
\end{example}

The scenario in the preceding example happens more generally,
as we can take note of the fact that the open sets in $U$ were
essentially defined with respect to the basis of open intervals.
\begin{theorem}[Egg yolk property]
    \label{thm:basiscriterion}
    Let $\mathcal T$ be a topology on $X$ and
    let $\mathcal B$ be some collection of sets in $\mathcal T$.
    Then, $\mathcal B$ is a basis for $\mathcal T$ if and only if
    for all open sets $U$ and all $x\in U$, there exists $B\in\mathcal B$
    such that $x\in B\subseteq U$.
\end{theorem}

\begin{proof}
    If $\mathcal B$ is a basis, then any open $U$ can be written as
    $\bigcup_{\alpha\in I}B_\alpha$ for some
    subcollection $\{B_\alpha\}_{\alpha\in I}$ of $\mathcal B$.
    Then, any $x\in U$ is contained in $B_\beta$ for some $\beta\in I$,
    so $x\in B_\beta\subseteq U$.

    Conversely, if $\mathcal B$ has the given property,
    then given open $U$, associate with each $x\in U$ a set $B_x\in\mathcal B$
    such that $x\in B_x\subseteq U$. Then,
    \begin{equation*}
	\bigcup_{x\in U}B_x = U
    \end{equation*}
    so $U$ is a union of sets in $\mathcal B$, meaning $\mathcal B$ is a basis
    for the given topology on $X$.
\end{proof}

\begin{remark}
    We called this the ``egg yolk property'' in class because
    the picture of $x\in U\subseteq X$, when drawing $U$ and $X$ as blobs,
    looks like an egg $X$ with a yolk $U$.
    This property is really a more general principle:
    if for each $x\in U$ one can find an open set $V$
    (often a basis element) such that $x\in V\subseteq U$,
    then $U$ must be an open set as well.
\end{remark}

The previous theorem can be interpreted as saying that
given basis $\mathcal B$ for topology $\mathcal T$ on $X$,
the sentence ``for all $x\in U$, there exists $B\in\mathcal B$ such that
$x\in B\subseteq U$'' means the same thing as ``$U$ is open''.
Observing this lets us quickly prove a nice criterion
for checking whether two topologies of $X$ are equal
just by looking at a basis of each topology.
\begin{theorem}[Bases give same topology]
    \label{thm:sametopologyfrombases}
    Let $\mathcal B_1$ be a basis for topology $\mathcal T_1$ on $X$, and
    let $\mathcal B_2$ be a basis for topology $\mathcal T_2$ on $X$.
    Then, $\mathcal T_1=\mathcal T_2$ if and only if for all $x\in X$,
    given $B_1\in\mathcal B_1$ containing $x$, there is a $B_2\in\mathcal B_2$
    such that $x\in B_2\subseteq B_1$, and vice versa.
\end{theorem}

\begin{proof}
    Observe that the given condition is equivalent to saying that
    each $B_1$ is open with respect to $\mathcal T_2$ and vice versa;
    therefore, if $U$ is open in $\mathcal T_1$,
    it is a union of sets in $\mathcal B_1$, hence also a union of sets
    that are open in $\mathcal T_2$, hence also open with respect to $\mathcal T_2$;
    and vice versa. That is, $\mathcal T_1$ and $\mathcal T_2$
    have the same collection of open sets!
\end{proof}

Thus, the topologies on $\RR^2$ generated by (1) the rectangles $(a,b)\times(c,d)$
and (2) the $\eps$-balls $B_\eps(x) = \{y\in\RR^2 : |x-y|<\eps \}$
are easily shown to be the same.

It is also sometimes helpful to consider a more minimal way of generating a topology:
\begin{definition}[Subbasis]
    \label{def:subbasis}
    Let $S$ be a collection of subsets of $X$ whose union is $X$.
    Then, the \vocab{topology generated by $S$}, denoted $\mathcal T(S)$,
    consists of arbitrary unions of finite intersections of sets in $S$.
    $S$ is called a \vocab{subbasis} of $\mathcal T(S)$.
\end{definition}

\begin{exercise}
    Check that $\mathcal T(S)$ is indeed a topology
    (thus the set of finite intersections of elements of $S$
    is a basis of $\mathcal T(S)$).
    Also, convince yourself that a basis for a topology
    is also a subbasis for the same topology.
\end{exercise}

It is helpful to think of the generated topology as
the ``smallest'' topology on $X$ that contains $S$, in the following sense:
\begin{definition}[Fine/coarse]
    \label{def:finecoarse}
    Let $\mathcal T_1$ and $\mathcal T_2$ be two topologies on a set $X$.
    We say that $\mathcal T_1$ is \vocab{finer} (or larger)
    than $\mathcal T_2$ and $\mathcal T_2$ is \vocab{coarser} (or smaller)
    than $\mathcal T_1$ if $\mathcal T_1\supseteq\mathcal T_1$.
\end{definition}
That is:
\begin{theorem}[Generated topology is minimal]
    \label{thm:minimality-generatedtop}
    Let $S$ be a subbasis of the topology $\mathcal T$.
    Then, $\mathcal T$ is the unique coarsest topology containing $S$,
    meaning if $\mathcal T'$ contains $S$ as well,
    then $\mathcal T\subseteq\mathcal T'$.
\end{theorem}

\begin{proof}
    Clearly if $\mathcal T'$ contains $S$ then it contains
    all finite intersections of elements of $S$ and
    arbitrary unions of these finite intersections,
    so $\mathcal T\subseteq\mathcal T'$.

    Furthermore, if $\mathcal T$ and $\mathcal T'$ both contain $S$
    and are as coarse as possible, then by definition of ``maximally coarse''
    we have both
    $\mathcal T\subseteq\mathcal T'$ and $\mathcal T'\subseteq\mathcal T$.
\end{proof}

\begin{exercise}
    Let $S$ be a subbasis for the topology $\mathcal T(S)$ on $X$.
    Show that $\mathcal T(S)$ is the intersection
    of all the topologies on $X$ containing $S$
    (you may want to first convince yourself that, in general,
    the intersection of two topologies on $X$ is also a topology on $X$).
\end{exercise}

\begin{exercise}
    Show that the set of open rays, in $\RR$, i.e.
    \begin{equation*}
	\{(a,\infty) : a\in\RR\}\cup\{(\infty,b) : b\in\RR\},
    \end{equation*}
    is a subbasis for the Euclidean topology on $\RR$.
\end{exercise}

\section{August 27}\label{sec:aug27}
\begin{definition}[Product topology]
    \label{def:finiteprodtop}
    Let $X$ and $Y$ be topological spaces.
    The \vocab{product topology} on $X\times Y$ is given by the basis
    \begin{equation*}
	\{U\times V : U\subseteq X\text{ and }V\subseteq Y\text{ are open}\}.
    \end{equation*}
    (In other words, open sets in $X\times Y$ look like
    arbitrary unions of boxes.)
\end{definition}

\begin{exercise}
    Verify that this basis really is a basis for $X\times Y$
    (as opposed to a subbasis; you'll need to check that the set is
    closed under finite intersections, see Definition \ref{def:subbasis}).
\end{exercise}

\begin{lemma}
    \label{lem:prodofbases}
    If $\mathcal B_X$ and $\mathcal B_Y$ are bases for
    topologies on $X$ and $Y$, respectively, then
    the product topology has basis
    $\mathcal B = \{U\times V : U\in\mathcal B_X,\,V\in\mathcal B_Y\}.$
\end{lemma}

\begin{proof}
    Let $(x,y)$ be a point in $X\times Y$ and $U$ an open set in $X\times Y$.
    Then, by definition of product topology, there exists open $U_X\subseteq X$
    and $U_Y\subseteq Y$ such that $(x,y)\in U_X\times U_Y\subseteq U$.
    Pick a basis element $B_X\in\mathcal B_X$ so that
    $x\in B_X\subseteq U_X$, and similarly pick $B_Y\in\mathcal B_Y$ so
    $y\in B_Y\subseteq U_Y$. Then we get
    \begin{equation*}
	(x,y)\in B_X\times B_Y\subseteq U_X\times U_Y\subseteq U
    \end{equation*}
    so the egg yolk property is satisfied and $\mathcal B$ is a basis.
\end{proof}

Thus, one can impose a topology on $\RR^2$ by giving open rectangles as a basis.
Of course, this is equal to the topology given by $\eps$-balls,
as shown in the case of $n=2$ last lecture.

\begin{definition}[Subspace topology]
    \label{def:subpsace}
    Let $Y$ be a subset of topological space $X$.
    Then, $Y$ can be made a topological space
    by declaring a subset $U$ of $Y$ to be open
    if and only if $U = V\cap Y$ for some open subset $V$ of $X$.
    This is called the \vocab{subspace topology}
    or \vocab{topology induced by $X$} or other phrases of the same ilk.
\end{definition}

\begin{example}
    Give $X = [0,1)\cup\{2\}$ the subspace topology it inherits
    as a subset of $\RR$.
    Then, $\{2\}$ and $[0,1)$ are open subsets of $X$,
    even though they aren't open in $\RR$
    (check the egg yolk property at $0$ and $2$ to show this).
\end{example}

\begin{exercise}
    Given a basis $\mathcal B_X$ for $X$ and $Y\subseteq X$,
    what is the obvious basis $\mathcal B_X$ for the topology of $Y$?
    Prove that this really is a basis.
\end{exercise}

\begin{definition}[Closed set]
    \label{def:closedset}
    A subset $E\subseteq X$ is closed if it is equal to $U^c$,
    the complement $X\setminus U$, of an open subset $U$.
\end{definition}

\begin{example}
    The set $[0,1]$ is closed in $\RR$ because it is the complement of
    the union of open sets $(-\infty,0)\cup(1,\infty)$,
    while $[0,1)$ is neither open nor closed --- it isn't open
    because egg yolk property fails at $0$,
    and it isn't closed because its complement fails egg yolk property at $1$.
\end{example}

\begin{prop}
    \label{prop:closedset-intersectionunion}
    Closed sets satisfy the following union and intersection properties:
    \begin{enumerate}[(1)]
	\ii $\varnothing$ and $X$ are closed.
        \ii Arbitrary intersections of closed sets are closed.
	\ii Finite unions of closed sets are closed.
    \end{enumerate}
    In principle, one could use closed sets to define a topology
    instead of open sets; the convention just happens to be open sets
    for most purposes (although for example the Zariski topology
    cares more about closed sets).
\end{prop}

\begin{proof}
    Immediate from writing closed sets as complements of open sets,
    since the complement passes through the unions and intersections.
\end{proof}

\section{September 3}\label{sec:sep03}

\subsection{Closure and limit points}

\begin{definition}[Interior and closure]
    \label{def:interior-closure}
    The \vocab{interior} $\accentset{\circ}{A}$ of $A\subseteq X$
    is the union of all open subsets of $X$ contained in $A$.
    The \vocab{closer} $\ol A$ of $A\subseteq X$
    is the intersection of all closed subsets of $X$ containing $A$.
\end{definition}

\begin{prop}
    \label{prop:interior-closure-extremal}
    The interior and closure are respectively maximal and minimal
    in the following sense:
    \begin{enumerate}[(1)]
        \ii $\interior A$ is the largest open set contained in $A$.
	\ii $\ol A$ is the smallest closed set containing $A$.
    \end{enumerate}
\end{prop}

This is of course immediate from the definition.
Note also that $A = \interior A$ is equivalent to $A$ being open
and $A = \ol A$ is equivalent to $A$ being closed.

The closure of a set can also be defined in terms of
more local properties of the set.
To make this idea more precise, we define the following:
\begin{definition}[Limit point]
    \label{def:limit-point}
    Let $A\subseteq X$. Then $x\in X$ is a \vocab{limit point} of $A$
    if any open set containing $x$ intersects $A$ at a point besides $x$,
    i.e. $U\cap A$ contains a point not equal to $x$
    for all open $U\ni x$.
\end{definition}

\begin{theorem}[Limit points provide closure]
    \label{thm:closure-by-limit-pts}
    Given $A\subseteq X$, let $A'$ be the set of limit points of $A$ in $X$.
    Then $\ol A = A\cup A'$.
\end{theorem}

\begin{proof}
    If $x\notin A\cup A'$, then there is a neighborhood $U\ni x$
    that does not intersect $A$. This means that $X\setminus U$
    is a closed set containing $A$ that doesn't contain $x$,
    so $x$ is also not in $\ol A$. Thus $\ol A\subseteq A\cup A'$.

    Conversely, if $F\supseteq A$ is closed and $x$ is not in $F$,
    then $x$ is contained in the open set $X\setminus F$,
    which does not intersect $A$ at all, so $x$ isn't a limit point for $A$.
    Therefore any closed set containing $A$ contains $A\cup A'$,
    so $\ol A$ contains $A\cup A'$.

    We have shown that $\ol A$ and $A\cup A'$ contain each other,
    so the two sets are equal.
\end{proof}

At this point a useful condition to impose on topological spaces is
that they be \vocab{Hausdorff}, meaning that
given any distinct $x$ and $y$ in $X$,
there exist open $U\ni x$ and $V\ni y$ that are disjoint.
This ensures that our space has enough open sets
to meaningfully separate points from each other;
most topological spaces we encounter in the real world are Hausdorff.

\begin{exercise}
    Show that $\RR^n$ with the standard topology is Hausdorff.
\end{exercise}

\begin{prop}
    \label{prop:finite-subset-of-hausdorff-closed}
    Finite subsets of Hausdorff spaces are closed.
\end{prop}

\begin{proof}
    It suffices to prove this for singleton sets,
    since finite subsets are finite unions of singleton sets.
    Indeed, if $x\in X$, then for any $y\in X$ not equal to $x$
    there are neighborhoods $U\ni x$ and $V\ne y$
    such that $U\cap V=\varnothing$ so in particular $x\notin V$;
    this implies that any $y\ne x$ is not a limit point of $\{x\}$,
    so $\{x\}$ is closed in $X$.
\end{proof}

\subsection{Continuous functions}
The standard definition of a continuous function between two metric spaces
asks that for any threshold $\eps$, there is a $\delta$ such that
two points that are at most $\delta$ apart in the domain
well be at most $\eps$ apart in the codomain.
Said differently, given $x\in X$, the ball $B_Y(f(x),\eps)$
has pre-image containing some open ball $B_X(x,\delta)$ in $X$.
By taking sub-balls of $B_Y(f(x),\eps)$ centered at each point in this ball,
we find that this forces $f\inv(B_Y(f(x),\eps))$ to be open in $X$
since we have an open ball in $X$ centered at each point of the pre-image
while still living inside the pre-image.
Because the balls in $Y$ are a basis for $Y$'s topology,
this shows that the pre-image of any open set in $Y$
must be open in $X$ (since unions pass through pre-images).

What we have done, in short, is eliminated any reference to
the metrics on $X$ and $Y$ in our definition of continuity,
referring instead to only topological data! This is the insight behind
the definition of continuous functions between topological spaces:
\begin{definition}[Continuous function, homeomorphism]
    \label{def:continuity-homeomorphism}
    A function $f\colon X\to Y$ between topological spaces
    is \vocab{continuous} if, for every open $V\subseteq Y$,
    the pre-image $f\inv(V)\subseteq X$ is open.
    If $f$ is bijective and $f\inv$ is continuous as well,
    then $f$ is a \vocab{homeomorphism} and
    $X$ and $Y$ are \vocab{homeomorphic}.
\end{definition}

\begin{exercise}
    Exhibit a homeomorphism between $\RR$ (with the standard topology)
    and any open interval in $\RR$ (with the subspace topology from $\RR$).
\end{exercise}

\begin{exercise}
    Show that $f\colon[0,2\pi)\to S^1\subseteq\RR^2$
    (where $S^1$ is the unit circle in $\RR^2$)
    defined as $f(t) = (\cos t, \sin t)$ is continuous and bijective
    but not a homeomorphism, because its inverse is not continuous.
\end{exercise}

\pagebreak

\section{September 8}\label{sec:sep08}

\subsection{Product topology}

Recall that given finitely many topological spaces
$X_1$, $X_2$, \dots, $X_n$ we can impose a topology on
$X_1\times X_2\times\cdots\times X_n$
by decreeing its basis to be
\begin{equation*}
    \{U_1\times\cdots\times U_n : U_i\overset{\text{open}}\subseteq X_i\}.
\end{equation*}

If instead we are given an arbitrary collection $\{X_\alpha\}_{\alpha\in A}$
of topological spaces, then one way to impose a topology on
$\prod_{\alpha\in A}X_\alpha$ is:
\begin{definition}[Box topology]
    \label{def:box-top}
    The \vocab{box topology} on the product $\prod_{\alpha\in A}X_{\alpha}$
    of topological space $\{X_\alpha\}_{\alpha\in A}$
    is given by the basis
    \begin{equation*}
	\left\{\prod_{\alpha\in A}U_{\alpha} :
	U_\alpha\overset{\text{open}}\subseteq X_\alpha\right\}
    \end{equation*}
\end{definition}

We will shortly see that this topology on the product is insufficient,
and we prefer the following one instead:
\begin{definition}[Product topology]
    \label{def:inf-product-top}
    The \vocab{product topology} on $\prod_{\alpha\in A} X_{\alpha}$
    is given by the subbasis
    \begin{equation*}
	\mathcal S = \bigcup_{\alpha\in A}
	\left\{\pi_\alpha\inv(U_\alpha) :
	U_{\alpha}\overset{\text{open}}\subseteq X_\alpha\right\},
    \end{equation*}
    where $\pi_\alpha$ is the projection onto $X_\alpha$.
\end{definition}

\begin{exercise}
    What is the basis of the product topology
    given by the subbasis $\mathcal S$?
    (i.e. what is the collection of
    finite intersections of sets in $\mathcal S$?)
\end{exercise}

\begin{exercise}
    Show that the box and product topologies agree
    for products of finitely many spaces.
\end{exercise}

In other words, the product topology is defined as the coarsest topology
that makes the $\pi_\alpha$ maps continuous.
The $\pi_\alpha$ are continuous because we are specifically saying that
the pre-image of any open set in $U_\alpha$ is open
in the product topology --- this is literally the definition of $\mathcal S$.
It shortly follows that this is also the coarsest topology
allowing the projections to be continuous.

\begin{theorem}[Continuity via component functions]
    \label{thm:f-cont-iff-components-cont}
    Suppose $\prod_{\alpha\in A} X_\alpha$ is given the product topology.
    Then, a map $f\colon Y\to\prod_{\alpha\in A}X_\alpha$ is continuous
    if and only if the components functions $f_\alpha\colon Y\to X_\alpha$,
    defined as $\pi_\alpha\circ f$, are all continuous.
\end{theorem}

\begin{proof}
    If $f$ is continuous then the result is immediate
    because compositions of continuous functions are continuous.
    (This hasn't been proven in these notes yet; go check it.)

    On the other hand, if the $f_\alpha$ are continuous,
    we can verify that $f$ is continuous by checking that $f\inv(U)$ is open
    where $U$ is an element of the subbasis $\mathcal S$
    that generates the product topology. Taking $U = \pi_\alpha\inv(U_\alpha)$,
    \begin{equation*}
        f\inv(U) = f\inv(\pi_\alpha\inv(U_\alpha)) =
	(\pi_\alpha\circ f)\inv(U_\alpha) = f_\alpha\inv(U_\alpha).\qedhere
    \end{equation*}
\end{proof}

The kicker is that this theorem fails for the box topology!
\begin{example}
    [Box topology is bad]
    Let $\RR^\omega$ be the product of countably many copies of $\RR$,
    endowed with the box topology.
    Consider $f\colon\RR\to\RR^\omega$ given by
    $f(t) = (t, t, t, \dots)$; this definitely should be continuous
    in any reasonable topology. However,
    \begin{equation*}
	U = \prod_{n=1}^\infty (-1/n, 1/n)
    \end{equation*}
    is an open set whose pre-image is $\bigcap_{n=1}^\infty (-1/n,1/n) = \{0\}$,
    which isn't open in $\RR$!
    (More abstractly, being able to take too many sets in the product
    defining $U$ forced an infinite intersection of open sets
    in the pre-image, which needn't be open.)
\end{example}

\subsection{Metric spaces}
Recall the definition of a metric space:
\begin{definition}[Metric space]
    \label{def:metric-space}
    A \vocab{metric space} is a set $X$ endowed with a \vocab{metric}
    $d\colon X\times X\to\RR_{\ge0}$ such that
    \begin{enumerate}[(1)]
        \ii $d(x,y) = 0$ if and only if $x=y$;
	\ii $d(x,y) = d(y,x)$ for all $x,y\in X$;
	\ii $d(x,y)\le d(x,y) + d(z,y)$ for all $x,y,z\in X$.
    \end{enumerate}
\end{definition}

\begin{example}
    $\RR^n$ with the distance function
    \begin{equation*}
	d(x,y) = \sqrt{\sum_{n=1}^\infty (x_i-y_i)^2}
    \end{equation*}
    is a metric space; the only nontrivial part of checking this
    is showing the triangle inequality. One way to do this is
    to anchor $z$ to the origin (since $d(x-z,y-z) = d(x,y)$)
    and then expand and use the Cauchy-Schwarz inequality.
\end{example}

Then, metric spaces can be given a topology,
as seen in analysis class:
\begin{definition}[Metric topology]
    \label{def:metric topology}
    The \vocab{metric topology} on metric space $(X,d)$
    has basis consisting of \vocab{metric balls}, i.e. the sets
    $B_\eps(x) = \{y\in X : d(x,y) < \eps\}$.
\end{definition}

\begin{exercise}
    Verify that this is a basis for a topology
    by showing that the intersection of two balls
    satisfies the egg yolk property
    (that is, the intersection of the balls must be open,
    hence a countable union of balls).
\end{exercise}

\begin{exercise}
    Show that $\rho\colon\RR^n\times\RR^n\to\RR_{\ge0}$
    defined by $\rho(x,y) = \max_i |x_i-y_i|$ is a metric on $\RR^n$
    and that the corresponding topology is also the Euclidean topology.
    What do the metric balls in this space look like?
\end{exercise}

\begin{theorem}[Equivalence of metrics]
    \label{thm:equivalent-metrics}
    Let $d$ and $d'$ be two metrics on $X$.
    Then the topologies generated by $d$ and $d'$ are the same
    if and only if there exists a constant $C\ge1$ such that
    \begin{equation*}
	B^{d'}_{C\inv\eps}(x)\subseteq B^d_\eps(x)
	\subseteq B^{d'}_{C\eps}(x)
    \end{equation*}
    for all $x\in X$ and $\eps>0$.
\end{theorem}

\begin{definition}[Metrizable]
    \label{def:metrizable}
    A topological space $X$ with topology $\mathcal T$ is \vocab{metrizle}
    if there is a metric on $X$ whose metric topology is $\mathcal T$.
\end{definition}

We will talk about metrizability Much Later.

\section{Septmeber 10}\label{sec:sep10}

\begin{prop}
    \label{prop:metric-spaces-hausdorff}
    Metric spaces are Hausdorff.
\end{prop}

\begin{proof}
    Let $x\ne y$ be in $X$. Then, $d(x,y) = d_0 > 0$.
    Take the open balls $B_{d_0/2}(x)\ni x$ and $B_{d_0/2}(y)\ni y$.
    These two don't intersect, since if $z$ were in their intersection,
    then by triangle inequality
    \begin{equation*}
	d_0 = \frac {d_0}2 + \frac {d_0}2 > d(x,z) + d(z,y) > d(x,y) = d_0,
    \end{equation*}
    which is a contradiction.
\end{proof}

\begin{theorem}[$\eps$-$\delta$ continuity]
    \label{thm:eps-delta-continuity-metric-space}
    A function $f\colon (X,d_X)\to(Y,d_Y)$ is continuous
    if and only if for each $x_1,x_2\in X$ and $\eps>0$,
    there exists $\delta>0$ (possibly depending on $x_1$, $x_2$)
    such that
    \begin{equation*}
        d_X(x_1,x_2)<\delta\implies d_Y(f(x_1),f(x_2)) < \eps.
    \end{equation*}
\end{theorem}

\begin{exercise}
    Show that the operations $+$, $-$, $\times$, $\div$
    are continuous as maps from $\RR^2$ to $\RR$.
\end{exercise}

\begin{definition}[Convergence]
    \label{def:convergence}
    Let $x_1,x_2,\ldots\in X$ be a sequence in topological space $X$.
    We say that $(x_n)_n$ \vocab{converges} to $x\in X$ if
    every open set containing $x$ contains cofinitely many terms
    of the sequence.
\end{definition}

\begin{prop}
    \label{prop:hausdorff-limits-unique}
    If $X$ is Hausdorff and a sequence $(x_n)_{n\in\NN}$ converges
    to both $x$ and $x'$ in $X$, then $x = x'$.
\end{prop}

\begin{proof}
    If $x\ne x'$ then we pick disjoint open sets containing $x$ and $x'$;
    it is not possible for both of these open sets to have
    cofinitely many points of the sequence.
\end{proof}

\begin{theorem}[Convergent sequences provide limit points]
    \label{thm:limit-points-by-seq}
    Let $A\subseteq X$. If $(x_n)_{n\in\NN}$ is a sequence of points in $A$
    with $x = \lim_{n\to\infty} x_n$, then $x\in\ol A$.
    If $X$ is a metric space, then the converse is true;
    that is, if $x\in\ol A$, then there is a sequence of points in $A$
    converging to $A$.
\end{theorem}

\begin{remark}
    More generally the converse holds in first countable spaces.
\end{remark}

\begin{proof}
    If $x\in A$ then we are done.
    Assuming $x\notin A$, the definition of convergence says that
    any neighborhood of $x$ contains points in the sequence $(x_n)$,
    which are in $A$; none of these $x_n$ are equal to $x$
    by assumption, so $x$ is a limit points for $A$, hence in $\ol A$.

    For the converse, we again are immediately done if $x\in A$
    because we can take the constant sequence.
    Otherwise, take successively smaller open balls
    $B_{1/n}(x)$ around $x$; since these are open and contain $x$,
    they intersect $A$, so take $x_n\in B_{1/n}(x)\cap A$.
    Since any open set containing $x$ contains some $\eps$-ball
    and for any $\eps>0$ there exists $n\in\NN$ such that $\frac 1n < \eps$,
    we get that any open set containing $x$ contains cofinitely many
    of the $x_n$.
\end{proof}

We can use sequences to phrase continuity if the spaces are nice enough,
just like we did with maps between metric spaces:
\begin{theorem}[Sequential continuity if domain is metric space]
    \label{thm:seq-cont-dom-metric-sp}
    Let $X$ be a metric space and $Y$ a map.
    Then $f$ is continuous if and only if for every sequence
    $x_n\to x$ in $X$ we have $f(x_n)\to f(x)$ in $Y$.
\end{theorem}

\begin{proof}
    \todo{todo}
\end{proof}

\begin{exercise}
    Show that $f\colon\RR^2\setminus\{(0,0)\}\to\RR$ defined by
    $f(x,y) = \frac{xy}{x^2+y^2}$ cannot be extended
    to a continuous function on $\RR^2$.
\end{exercise}

\begin{definition}[Uniform convergence]
    \label{def:uniform-convergence}
    Let $f_n\colon X\to(Y,d_Y)$ be a sequence of functions
    (where $X$ can be anything it wants to be).
    They \vocab{converge uniformly} to $f\colon X\to Y$ if
    given $\eps>0$ there exists $N$ such that $d(f_n(x), f(x)) < \eps$
    for all $n\ge N$ and $x\in X$.
\end{definition}
In other words, $f_n(x)\to f(x)$ at roughly the same rate for all $x\in X$.

\begin{exercise}
    Let $f_n\colon[0,\half]\to\RR$ be defined by $f_n(x) = x^n$.
    Show that $f_n$ converges to the zero function uniformly.
    What if the domain is instead $[0,1]$ or $[0,1)$?
\end{exercise}

\section{September 15}\label{sec:sep15}
\subsection{Quotient topology and quotient maps}

Another way to construct topological spaces is by gluing pieces of
spaces together, such as rolling up a square into a cylinder
and then rolling that cylinder into a donut,
or simply rolling an interval into a circle.
One can make this idea precise as follows:

\begin{definition}[Quotient topology]
    \label{def:quotient-top}
    Let $X$ be a topological space endowed with an equivalence relation $\sim$.
    Then, the set of equivalence classes $X/{\sim}$ is the
    \vocab{quotient space} of $X$ modulo $\sim$,
    which is topologized by declaring $U\subseteq X/{\sim}$ to be open
    if and only if the union of the equivalence classes in $U$,
    i.e. $\bigcup_{[x]\in U} [x]$, is open in $X$.
    Equivalently, $U\subseteq X/{\sim}$ is open if and only if
    $p\inv(U)\subseteq X$ is open, where the projection $p\colon X\to X/{\sim}$
    is the map $p(x) = [x]$.
\end{definition}
Note of course that this means $p$ is a continuous map $X\to X/{\sim}$.

\begin{example}
    [$S^1$ is a quotient of $\lbrack0,2\pi\rbrack$]
    Give $[0,2\pi]$ the equivalence relation
    whose only nontrivial relation is $0\sim 2\pi$.
    Then $X=[0,2\pi]/{\sim}$ is homeomorphic to the unit circle
    $S^1 = \{(\cos\theta,\sin\theta) \in\RR^2 : \theta\in[0,2\pi)\}\subseteq\RR^2$
    via the map $h\colon X\to S^1$ defined as $h(x) = (\cos x,\sin x)$.
\end{example}

\begin{exercise}
    Check that $h$ really is a homeomorphism $X\to S^1$.
\end{exercise}

\begin{example}
    Let $X = \ol{B_1(0)}\subseteq\RR^2$ be the closed disk of radius $1$
    centered at $0$, and let $\sim$ be the relation such that
    any two points on the boundary $S^1$ are equivalent.
    This is actually homeomorphic to the sphere $S^2\subseteq\RR^3$!
    We can imagine the boundary collapsing onto the north pole,
    and then it is not too hard to convince yourself
    that open sets around the north pole in $S^2$ correspond to
    open sets in $X$.
\end{example}

\begin{exercise}
    Let $X = [0,1]^2$ and identify points with the relation
    $(x,0)\sim(x,1)$ for all $x\in[0,1]$ and $(0,y)\sim(1,0)$
    for all $y\in[0,1]$. What shape does $X/{\sim}$ look like?
    (Hint: identify the two pairs of edges one at a time.)
\end{exercise}

We can also get quotient maps in a sort of backwards way,
by deriving the underlying equivalence relation via
a map that behaves like the projection map.
\begin{definition}[Quotient map]
    \label{def:quotient-map}
    Let $p\colon X\to Y$ be a surjective map such that $p\inv(U)\subseteq X$
    is open if and only if $U\subseteq Y$ is open.
\end{definition}
In this case $Y$ can be thought of as the quotient of $X$ by the relation
$x\sim x'$ whenever $x,x'\in p\inv(Y)$ for some $y\in Y$.
\begin{exercise}
    Show that the map $p\colon[0,2\pi]\to S^1$ given by
    $p(t) = (\cos t,\sin t)$ is continuous and surjective
    but is not a quotient map.
\end{exercise}

\begin{prop}
    \label{prop:quotient-universal-property}
    Let $X/{\sim}$ be a quotient space of $X$ with
    projection map $p\colon X\to X/{\sim}$
    and suppose that there is a continuous map $f\colon X\to Z$
    that is constant on $p\inv(y)$ for each $y\in X/{\sim}$.
    Then, there exists a unique continuous map $\tilde f\colon X/{\sim}\to Z$
    such that $\tilde f\circ p = f$. 
\end{prop}

\begin{proof}
    For a given $y\in X/{\sim}$, define $\tilde f(y)$ to be
    the value of $f$ at any point in $p\inv(y)$.
    It is easy to see that $\tilde f\circ p = f$,
    so it remains to check continuity.
    \todo{check continuity}
\end{proof}

\begin{example}
    [Real projective space]
    $\RR P^n = S^n/{\sim}$ where $\sim$ identifies antipodal points.
    (TODO elaborate)
\end{example}

Given $f\colon S^1\to\RR$ defined by $f(x,y) = y^2$ if $y\ge0$
and $f(x,y)=0$ if $y<0$, does there exist a corresponding $\tilde f$?

\subsection{Connectedness}

\begin{definition}[Connectedness]
    \label{def:connectned}
    Let $X$ be a topological space with $X = U\cup V$
    where $U,V$ are nonempty disjoint open sets; we say that $\{U,V\}$ is
    a \vocab{separation} of $X$ and that $X$ is \vocab{connected}
    if it does not admit a separation.
\end{definition}

\todo{examples}

\section{September 17}\label{sec:sep17}

\begin{prop}
    \label{prop:connected-iff-no-clopen}
    $X$ is connected if and only if the only subsets of $X$
    that are both open and closed are $X$ and $\varnothing$.
\end{prop}

\begin{proof}
    There exists a set $U$ that is both open and closed
    if and only if there is a separation of $X$
    (which is equivalent to $X$ being disconnected),
    since $X = U\cup V$ is the union of two disjoint nonempty open sets
    if and only if $U = X\setminus V$ is open by definition
    and closed by being the complement of $V$.
\end{proof}

\begin{theorem}[Unions of connected sets]
    \label{thm:union-connected-sets}
    Let $\{U_\alpha\}$ be a collection of connected sets in $X$.
    If the intersection of the $U_\alpha$ is nonempty, then
    $\bigcup_\alpha U_\alpha$ is connected.
\end{theorem}

\begin{soln}
    Suppose for the sake of contradiction that $\bigcup_\alpha U_\alpha$
    has a separation $V\sqcup V'$.

    \todo{lazy}
\end{soln}

\begin{theorem}[$f(\text{connected}) = \text{connected})$]
    \label{thm:connected-pass-thru-cts}
    Let $f\colon X\to Y$ be continuous.
    If $X$ is connected then $f(X)$ is connected.
\end{theorem}

\begin{proof}
    Suppose the subspace $f(X)\subseteq Y$ has a separation $U\sqcup V$.
    Then $U = \tilde U \cap f(X)$ and $V = \tilde V\cap f(X)$
    for some $\tilde U$, $\tilde V$ open in $Y$.
    Then the pullbacks of $\tilde U$ and $\tilde V$ are
    clearly open (by continuity) and nonempty, disjoint, and cover $X$
    (by basic properties of pullbacks of sets),
    so $f\inv(\tilde U)$ and $f\inv(\tilde V)$ are a separation of $X$,
    hence $X$ is disconnected if $f(X)$ is disconnected.
\end{proof}

\begin{definition}[Path connected]
    \label{def:path-connected}
    $X$ is \vocab{path connected} if there exists a continuous path
    between any two points $x,y\in X$; that is, there is a continuous function
    $f\colon[0,1]\to X$ such that $f(0)=x$ and $f(1)=y$.
\end{definition}

\begin{prop}
    \label{prop:path-conn-impiles-conn}
    
\end{prop}
\todo{copy from 320 notse}

\todo{comb space (line segments and $(0,1)$)}

\section{September 22}\label{sec:sep22}

\begin{lemma}
    \label{lem:set-closure-sandwich-connected}
    Let $X$ be a topological space and let $A\subseteq X$ be connected.
    If $A\subseteq B\subseteq\ol A$ then $B$ is connected.
\end{lemma}

\begin{corollary}
    \label{coro:deleted-comb-space-connected}
    The deleted comb space is connected.
\end{corollary}

\begin{proof}[Proof of \autoref{lem:set-closure-sandwich-connected}]
    Assume $B$ is disconnected, so there exist nonempty sets
    $U,V\subseteq B$ that are open in $B$ whose union is $B$.
    Let $\widetilde U = U\cap A$ and $\widetilde V = V\cap A$.
    Then $\widetilde U$ and $\widetilde V$ are open in $A$
    and $A = \widetilde U\cup\widetilde V$, so it remains to show that
    $\widetilde U$, $\widetilde V$ are nonempty.
    By definition of the subspace topology there exists
    $\widehat U$ that is open in $X$ such that $U = \widehat U\cap B$,
    hence $\widetilde U = \widehat U\cap A$. By the closure condition
    either some point of $A$ is in $\widehat U$, in which case we're done,
    or $\widehat U$ contains a limit point of $A$, in which case
    it must still contain a point of $A$, so we are still done
    (because either way $\widetilde U$ is nonempty,
    hence so is $\widetilde V$ by symmetry).
\end{proof}

\begin{theorem}[Intervals are connected]
    \label{thm:intervals-are-connected}
    Intervals and rays in $\RR$
    (i.e. sets of the form $(a,b)$, $(a,b]$, $[a,b)$, or $[a,b]$
    such that $-\infty\le a\le b\le\infty$, with strict inequalities
    if the corresponding endpoint is closed)
    are connected.
\end{theorem}

\begin{proof}
    Let $I$ be an interval and $U\cup V$ is a separation.
    blah blah\todo{finish}
\end{proof}

\begin{corollary}
    [Intermediate value theorem]
    \label{coro:ivt}
    blah blah
\end{corollary}

\begin{corollary}
    [Brouwer fixed point theorem, 1D]
    \label{coro:brouwer-1d}
\end{corollary}

\begin{exercise}
    Does the equation $x^3 + x = e^{\sin(x^3)} + 10x^2 - \log(\abs x+1)$
    have a solution in $\RR$?
\end{exercise}

\section{September 24}\label{sec:sep24}
Given a topological space, one way to look at
the connectedness/path-connectedness properties
can be understood by looking at maximal (path-)connected subsets.

\begin{definition}[(Path) connected components]
    \label{def:connected-components}
    A \vocab{connected component} (or just \vocab{component})
    of a space $X$ is a maximal connected subset
    of $X$. Equivalently, it is an equivalence class of $X$
    under the relation $a\sim b$ if there exists a connected subset of $X$
    containing both $a$ and $b$.
    Analogously, \vocab{path components} are connected components
    but for path connectedness.
\end{definition}

It is easy to show (and you should convince yourself if it isn't obvious)
that any (path) connected subset of $X$ can only lie in one component of $X$.

\begin{definition}[Local (path) connectedness]
    \label{def:local-connectedness}
    We say that $X$ is locally (path) connected if
    for any $x\in X$ and neighborhood $U\ni x$
    there exists an open (path) connected $V\subseteq U$ such that $V\ni x$.
\end{definition}

\begin{example}
    Local connectedness and connectedness do not imply each other.
    \begin{itemize}
        \ii $\RR$ (or any connected manifold) is locally connected
	and locally connected.
	\ii $[0,1)\cup(1,2]$ is locally connected
	but clearly not connected.
	\ii The comb space is connected but not locally connected
	(find a neighborhood that breaks!).
	\ii $\QQ$ is neither connected nor locally connected.
    \end{itemize}
\end{example}

\begin{theorem}[Locally connected means components of open are open]
    \label{thm:lc-iff-components-of-open-are-open}
    $X$ is locally connected if and only if each component of
    any open $U\subseteq X$ is open.
\end{theorem}

\begin{theorem}[Path component properties]
    \label{thm:path-component-props}
    \begin{enumerate}[(1)]
        \ii Path components lie inside connected components.
	\ii If $X$ is locally path connected then the components
	and path components of $X$ are the same.
    \end{enumerate}
\end{theorem}










\end{document}
